\section{Discussion: what failure of admissibility means}
\subsection{Storage coordinates need not be terminal-current coordinates}
The real-data study makes a distinction necessary: a conservative magnetic
subsystem does not automatically imply a potential whose gradient against
measured terminal current is the reconstructed flux. A schematic loss branch
has $\statorcurrent=\magnetizingcurrent+\vect i_{\rm Fe}$.
The storage relation can be conservative in $\magnetizingcurrent$ while its
stationary reduction against $\statorcurrent$ is not. Richter et al.
\cite{richter2014} distinguish magnetic-state matching from equality of
terminal currents in motoring/generating measurements.

For an algebraic illustration, take
$\flux=\matr L\magnetizingcurrent$ with
$\matr L=\operatorname{diag}(L_d,L_q)$ and
$\vect i_{\rm Fe}=c\rotationgenerator\flux$, $c=\omega_e/R_{\rm Fe}$.
Its dissipation in the same dq normalization is
$\|\omega_e\rotationgenerator\flux\|^2/R_{\rm Fe}>0$ for nonzero branch voltage;
physical three-phase power adds the $3/2$ factor under amplitude-invariant
coordinates. The storage law itself remains conservative. Eliminating its branch
current yields
\begin{equation}
 \frac{\partial\flux}{\partial\vect i_s}
 =\matr L(\matr I+c\rotationgenerator\matr L)^{-1}
 =\frac{1}{1+c^2L_dL_q}
 \begin{bmatrix}L_d&cL_dL_q\\-cL_dL_q&L_q\end{bmatrix}.
\end{equation}
Its terminal-current curl is $-2cL_dL_q/(1+c^2L_dL_q)$, despite conservative
storage magnetics. This illustrative constant branch is not identified for
the supplied machine. Indeed, its antisymmetric affine cross terms give
equal symmetry and curl equivalents; it does not automatically explain their
observed difference. Iron loss remains one open model component, not a
demonstrated cause.

Nor does simply relabeling currents transform a differential form correctly.
If $\magnetizingcurrent=f(\statorcurrent)$, then
\begin{equation}
 \nabla_{\vect i_s}\bigl[W'_{\rm mag}(f(\vect i_s))\bigr]
 =Df(\vect i_s)\trans\flux_{\rm mag}(f(\vect i_s)).
 \label{eq:current-pullback}
\end{equation}
The Jacobian transpose is required. This gradient is generally not the
unchanged flux vector plotted against new current labels. Magnetic coordinates,
their conjugate quantities, fixed state, and the subsystem being fitted must
therefore be specified together. Loss-aware synchronous-machine context is
distinct from induction-machine results; the Kullick--Hackl journal study
\cite{kullick2023} concerns induction machines and supplies no PMSM
co-energy validation for these terminal coordinates.
